% Part: normal-modal-logic
% Chapter: syntax-and-semantics
% Section: schemas

\documentclass[../../../include/open-logic-section]{subfiles}

\begin{document}

\olfileid{nml}{syn}{sch}

\olsection{పథకాలు మరియు చెల్లుబాటుతనం}

\begin{defn}
  \emph{పథకం} అంటే ఏదో మోడల్
  \tetoken{సూత్రం}{!!{formula}}~$!C$ యొక్క
  ప్రతిస్థాపన నిదర్శనాలన్నీ, అవి మాత్రమే గల
  \tetoken{సూత్రాల}{!!{formula}s} సమితి; అంటే,
  \[
  \Setabs{!B}{\lexists[!D_1], \dots, \lexists[!D_n] \left(!B =
  \SSubst{!C}{\subst{!D_1}{p_1}, \dots, \subst{!D_n}{p_n}} \right) }.
  \]
  \tetoken{సూత్రం}{!!{formula}}~$!C$ను ఆ పథకపు
  \emph{లక్షణ} \tetoken{సూత్రం}{!!{formula}} అంటాం.
  \tetoken{ప్రతిజ్ఞావాక్య చరాల}{!!{propositional
  variable}s} పేర్ల మార్పు వరకు అది ఏకైకం.
  \tetoken{సూత్రం}{!!^a{formula}}~$!A$ ఆ సమితిలో
  సభ్యమైతే, అది పథకానికి \emph{నిదర్శనం}.
\end{defn}

$!C$ పరమాణు భాగాల స్థానంలో ‘$!A$’, ‘$!B$’,
~\dots ఉంచి వచ్చిన అధిభాషా వ్యక్తీకరణతో
పథకాన్ని సూచించడం సౌకర్యకరం. ఉదాహరణకు,
‘$!A$’, ‘$!A \lif \Box!A$’,
‘$!A \lif (!B \lif !A)$’ పథకాలను సూచిస్తాయి.
వాటి లక్షణ \tetoken{సూత్రాలు}{!!{formula}s}
వరుసగా $p$, $p \lif \Box p$, $p \lif (q
\lif p)$. ‘$!A$’ పథకం \emph{అన్ని}
\tetoken{సూత్రాల}{!!{formula}s} సమితిని సూచిస్తుంది.

\begin{defn}
  ఒక పథకం నమూనాలో \emph{సత్యం} అంటే దాని
  నిదర్శనాలన్నీ ఆ నమూనాలో సత్యం;
  ఒక పథకం \emph{చెల్లుబాటు} అంటే అది
  ప్రతి నమూనాలో సత్యం.
\end{defn}

\begin{prop}\ollabel{prop:Kvalid}
  కింది \Ax{K} పథకం చెల్లుబాటవుతుంది:
  \begin{equation*}
  \Box(!A \lif !B) \lif (\Box !A \lif \Box !B). \tag{\Ax{K}}
  \end{equation*}
\end{prop}

\begin{proof}
  పథకంలోని ప్రతి నిదర్శనం ప్రతి నమూనాలోని
  ప్రతి లోకంలో సత్యమని చూపాలి. ఏదైనా
  $\mModel{M} = \tuple{W,R,V}$ను,
  $w \in W$ను తీసుకుందాం. ఒక లోకంలో
  సోపాధికం సత్యమని చూపడానికి, దాని
  పూర్వాంగం సత్యమనుకొని ఉత్తరాంగమూ
  సత్యమని చూపుతాం. ఇక్కడ
  $\mSat{M}{\Box(!A \lif !B)}[w]$,
  $\mSat{M}{\Box !A}[w]$ అనుకుందాం.
  $\mSat{M}{\Box !B}[w]$ చూపాలి.
  $Rww'$ అయ్యే ఏదైనా $w'$ను తీసుకుందాం.
  మొదటి ఊహ ప్రకారం
  $\mSat{M}{!A \lif !B}[w']$;
  రెండవ ఊహ ప్రకారం $\mSat{M}{!A}[w']$.
  అందువల్ల $\mSat{M}{!B}[w']$.
  $w'$ ఏదైనా కాబట్టి $\mSat{M}{\Box !B}[w]$.
\end{proof}

\begin{prop}\ollabel{prop:Dual-valid}
  కింది \Dual{} పథకం చెల్లుబాటవుతుంది:
  \begin{equation*}
  \Diamond !A \liff \lnot\Box\lnot !A. \tag{\Dual}
  \end{equation*}
\end{prop}

\begin{proof}
  సాధన.
\end{proof}

\begin{prob}
  \olref[nml][syn][sch]{prop:Dual-valid}ను నిరూపించండి.
\end{prob}

\begin{prop}\ollabel{prop:soundMP}
  ఒక నమూనాలోని ఒక లోకంలో $!A$,
  $!A \lif !B$ సత్యమైతే $!B$ కూడా
  అక్కడ సత్యం. కాబట్టి చెల్లుబాటయ్యే
  \tetoken{సూత్రాల}{!!{formula}s} సమితి
  మోడస్ పోనెన్స్ కింద సంవృతం.
\end{prop}

\begin{prop}\ollabel{prop:valid-instances}
  \tetoken{సూత్రం}{!!^a{formula}}~$!A$
  చెల్లుబాటవుతుంది, అప్పుడే, అప్పుడే
  దాని ప్రతిస్థాపన నిదర్శనాలన్నీ
  చెల్లుబాటవుతాయి. మరోలా అంటే,
  పథకం చెల్లుబాటవుతుంది, అప్పుడే,
  అప్పుడే దాని లక్షణ
  \tetoken{సూత్రం}{!!{formula}} చెల్లుబాటవుతుంది.
\end{prop}

\begin{proof}
  “అయితే” దిశ స్పష్టం: $!A$
  తనకే ఒక ప్రతిస్థాపన నిదర్శనం.

  “అప్పుడే” దిశకు ఈ వాదాన్ని చూపుదాం:
  $\mModel{M} = \tuple{W, R, V}$
  మోడల్ నమూనా అనుకుందాం; $!B \ident
  \SSubst{!A}{\subst{!D_1}{p_1},\dots,\subst{!D_n}{p_n}}$
  అనేది $!A$కి ప్రతిస్థాపన నిదర్శనం.
  $V'(p_i) = \Setabs{w}{\mSat{M}{!D_i}[w]}$
  అయ్యేలా $\mModel{M'} = \tuple{W, R,
    V'}$ను నిర్వచిద్దాం; ఇతర చరాలకు
  $V'(q)=V(q)$.
  \sourcecorrection{OLTENMLSCH-001}{మూలంలో
    $V'$ను $p_i$లకే నిర్దేశించారు.
    మిగతా చరాలకు అసలు కేటాయింపునే
    ఉంచి నమూనా నిర్వచనాన్ని పూర్తిచేశాం.}
  ఏ $w \in W$కైనా $\mSat{M}{!B}[w]$
  అప్పుడే, అప్పుడే $\mSat{M'}{!A}[w]$.
  (దీని నిరూపణను సాధనగా విడిచిపెడుతున్నాం.)
  ఇప్పుడు $!A$ చెల్లుబాటై,
  $!A$కి చెందిన ఏదో ప్రతిస్థాపన నిదర్శనం
  $!B$ చెల్లుబాటు కాలేదనుకుందాం.
  అప్పుడు ఏదో $\mModel{M} = \tuple{W, R, V}$
  మరియు ఏదో $w \in W$కు
  $\mSat/{M}{!B}[w]$. కానీ వాదం ప్రకారం
  $\mSat/{M'}{!A}[w]$; కనుక $!A$
  చెల్లుబాటు కాదు. ఇది వైరుధ్యం.
\end{proof}

\begin{prob}
  \olref[nml][syn][sch]{prop:valid-instances}
  నిరూపణలోని “అప్పుడే” భాగంలో
  చెప్పిన వాదాన్ని నిరూపించండి.
  (సూచన: $!A$పై ఆగమనం ఉపయోగించండి.)
\end{prob}

అయితే, పథకం నమూనాలో సత్యమైతే
అప్పుడే దాని లక్షణ సూత్రం ఆ
నమూనాలో సత్యమని చెప్పలేం.
“అప్పుడే” దిశ మాత్రం నిలుస్తుంది:
$!A$లోని ప్రతి నిదర్శనం
$\mModel{M}$లో సత్యమైతే
$!A$~కూడా $\mModel{M}$లో సత్యమే.
కానీ $!A$ అనేది $\mModel{M}$లో
సత్యమై ఉండి, $!A$కి చెందిన ఏదో నిదర్శనం
$\mModel{M}$లోని ఏదో లోకంలో
అసత్యమయ్యే అవకాశం ఉంది. సరళ
ప్రతిదృష్టాంతంగా, ఒక్క లోకం~$w$
మాత్రమే ఉండి $V(p) = \{w\}$
అయ్యే నమూనాలో~$p$ను పరిశీలించండి.
అప్పుడు $p$ అనేది~$w$లో సత్యం.
కానీ $\lfalse$ అనేది~$p$కి
నిదర్శనం; అది~$w$లో సత్యం కాదు.

\begin{prob}
  కింది \tetoken{సూత్రాలలో}{!!{formula}s}
  ఏదీ చెల్లుబాటు కాదని చూపండి:
  \begin{enumerate}
    \item[\Ax{D}:] \quad $\Box p \lif \Diamond p$;
    \item[\Ax{T}:] \quad $\Box p \lif p$;
    \item[\Ax{B}:] \quad $p \lif \Box\Diamond p$;
    \item[\Ax{4}:] \quad $\Box p \lif \Box \Box p$;
    \item[\Ax{5}:] \quad $\Diamond p \lif \Box \Diamond
      p$.
  \end{enumerate}
\end{prob}

\begin{table}[t]
    \centering
    \begin{tabular}{| l || l |}
      \hline
      {\emph{చెల్లుబాటయ్యే పథకాలు}} & {\emph{చెల్లుబాటు కాని పథకాలు}} \\
      \hline\hline
      $\Box(!A \lif !B) \lif (\Diamond !A \lif \Diamond !B)$
      & $\Box (!A \lor !B) \lif (\Box !A \lor \Box !B)$ \\
      $\Diamond (!A \lif !B) \lif (\Box !A \lif \Diamond
      !B)$
      & $(\Diamond !A \land \Diamond !B) \lif \Diamond (!A
      \land !B)$\\
      $\Box (!A \land !B) \liff (\Box !A \land \Box !B)$
      & $!A \lif \Box !A$ \\
      $\Box !A \lif \Box (!B \lif !A)$
      & $\Box \Diamond !A \lif !B$ \\
      $\lnot \Diamond !A \lif \Box (!A \lif !B)$
      & $\Box \Box !A \lif \Box !A$ \\
      $\Diamond (!A \lor !B) \liff (\Diamond !A \lor
      \Diamond !B)$
      & $\Box \Diamond !A \lif \Diamond \Box !A$. \\
      \hline
    \end{tabular}
    \caption{చెల్లుబాటయ్యే మరియు చెల్లుబాటు కాని పథకాలు.}
    \ollabel{tab:valid-invalidSchemas}
\end{table}

\begin{prob}%\ollabel{ex:in/validSchemas}
  \olref[nml][syn][sch]{tab:valid-invalidSchemas}లో
  మొదటి నిలువు వరుస పథకాలు
  చెల్లుబాటవుతాయని, రెండవ నిలువు వరుస
  పథకాలు చెల్లుబాటు కావని నిరూపించండి.
\end{prob}

\begin{prob}
  కింది పథకాలు చెల్లుబాటవుతాయా,
  కావా నిర్ణయించండి:
  \begin{enumerate}
  \item $(\Diamond !A \lif \Box !B) \lif (\Box !A \lif \Box
    !B)$;
  \item $\Diamond(!A \lif !B) \lor \Box(!B \lif !A)$.
  \end{enumerate}
\end{prob}

\begin{prob}
  కింది ప్రతి పథకానికి, ఆ
  \tetoken{సూత్రం}{!!{formula}}
  యొక్క ప్రతి నిదర్శనం
  $\mModel{M}$లో సత్యమయ్యేలా
  ఒక నమూనా $\mModel{M}$ను కనుగొనండి:
  \begin{enumerate}
  \item $p \lif \Diamond\Diamond p$;
  \item $\Diamond p \lif \Box p$.
  \end{enumerate}
\end{prob}

\end{document}
